\chapter{Introduction}
\label{sec:introduction}

Legal documents pose a distinctive challenge for natural language processing.
They are long, densely cross-referenced and written in a register that differs
considerably from the general-domain corpora on which most pre-trained models
are trained. Early neural approaches to legal judgment prediction showed that
general architectures transfer only partially to this setting
\parencite{chalkidis2019}.

The introduction of the Transformer \parencite{vaswani2017} and of large
pre-trained encoders such as \model{BERT} \parencite{devlin2019} changed this
picture. \textcite{chalkidis2020} demonstrated that further pre-training on
legal corpora yields consistent gains, and \textcite{zheng2021} later examined
the conditions under which such domain adaptation is worthwhile. As shown in
\cref{fig:pipeline}, this thesis follows the same general recipe.

\begin{figure}[t]
  \centering
  \includegraphics[width=0.8\textwidth]{example-image}
  \caption{Overview of the domain-adaptation pipeline used in this thesis.}
  \label{fig:pipeline}
\end{figure}

\section{Research Questions}
\label{sec:research-questions}

This work addresses whether domain-specific pre-training improves performance on
downstream legal classification tasks, and at what computational cost. The
experimental setup is described in \cref{sec:methodology}.
